\documentclass[12pt]{article}

\usepackage{amsmath}
\usepackage{amssymb}
\usepackage{graphicx}
\usepackage{hyperref}
\usepackage[margin=1in]{geometry}
\usepackage{natbib}
\bibliographystyle{aasjournal}

\newcommand{\bq}{\mathbf{q}}
\newcommand{\bx}{\mathbf{x}}
\newcommand{\bk}{\mathbf{k}}
\newcommand{\bPsi}{\boldsymbol{\Psi}}
\newcommand{\bv}{\mathbf{v}}
\newcommand{\bn}{\hat{\mathbf{n}}}
\newcommand{\Dp}{D_+}
\newcommand{\Dpp}{D_+^{(2)}}
\newcommand{\Omm}{\Omega_{\rm m}}
\newcommand{\calH}{\mathcal{H}}

\title{Lagrangian Perturbation Theory: Zel'dovich Approximation and 2LPT}
\author{Notes on IC generation with MUSIC2}
\date{\today}

\begin{document}

\maketitle
\tableofcontents
\newpage

%----------------------------------------------------------------------
\section{Introduction}

Cosmological $N$-body initial conditions (ICs) are generated by
displacing and accelerating particles from a regular lattice according
to a perturbed density field.  The standard framework is
\emph{Lagrangian perturbation theory} (LPT), in which one tracks
fluid elements by their initial (Lagrangian) positions $\bq$ and
solves for the displacement $\bPsi(\bq, t)$ that carries each element
to its Eulerian position $\bx$ at time $t$.

At first order this yields the \emph{Zel'dovich approximation} (ZA;
\citealt{Zeldovich1970}), which is exact for planar collapse and is
used for IC codes at high redshift.  Second-order LPT (2LPT;
\citealt{Buchert1993,Bouchet1995,Crocce2006}) adds a correction that
accounts for mode coupling and reduces transient artefacts caused by
initialising simulations at finite redshift.  \textsc{music2}
\citep{Hahn2011} implements both, controlled by the conf key
\texttt{use\_2LPT = yes}.

%----------------------------------------------------------------------
\section{Fluid Equations in an Expanding Universe}
\label{sec:fluid}

Before deriving the LPT displacement fields it is useful to record the
governing fluid equations in a $\Lambda$CDM background and their
linearised forms.

\subsection{Continuity, Euler, and Poisson equations}

In proper (physical) coordinates $\mathbf{r}$ the equations for a
self-gravitating pressureless fluid are
%
\begin{align}
  \frac{\partial\rho}{\partial t}
    + \nabla_r\cdot(\rho\bv) &= 0,
  \label{eq:continuity_proper}\\
  \frac{d\bv}{dt} &= -\nabla_r\psi,
  \label{eq:euler_proper}\\
  \nabla_r^2\psi &= 4\pi G\rho,
  \label{eq:poisson_proper}
\end{align}
%
where $d/dt = \partial/\partial t + \bv\cdot\nabla_r$ is the
convective derivative, $\psi$ is the Newtonian gravitational potential,
and $\nabla_r$ denotes the gradient in proper coordinates.

Switching to comoving coordinates $\bx = \mathbf{r}/a$ and splitting
each field into a homogeneous background plus a small perturbation,
$\bv = \bv_0 + \delta\bv$,
$\rho = \rho_0 + \delta\rho$,
$\psi = \psi_0 + \Phi$,
with Hubble flow $\bv_0 = H\mathbf{r}$ and $\Phi \equiv \delta\psi$
the perturbed gravitational potential, the linearised equations become
%
\begin{align}
  \dot\delta
    &= -\nabla\cdot\delta\bv,
  \label{eq:continuity_lin}\\
  \frac{d\,\delta\bv}{dt} + H\,\delta\bv
    &= -\frac{1}{a}\nabla\Phi,
  \label{eq:euler_lin}\\
  \nabla^2\Phi
    &= 4\pi G\rho_0\,a^2\,\delta
     = \tfrac{3}{2}H^2\Omm(a)\,a^2\,\delta,
  \label{eq:poisson_lin}
\end{align}
%
where $\nabla$ now denotes the gradient in comoving coordinates, dots
are derivatives with respect to cosmic time $t$, and we have used the
Friedmann equation $4\pi G\rho_0 = \tfrac{3}{2}H^2\Omm(a)$.

\subsection{Density contrast and linear growth equation}
\label{sec:growth}

The density contrast is defined as
%
\begin{equation}
  \delta(\bx, t) \equiv \frac{\delta\rho(\bx,t)}{\rho_0(t)}.
  \label{eq:delta_def}
\end{equation}
%
Combining Eqs.~(\ref{eq:continuity_lin})--(\ref{eq:poisson_lin}) gives
the second-order ODE for $\delta$ in a pressureless fluid:
%
\begin{equation}
  \ddot{\delta} + 2H\dot{\delta}
    = \tfrac{3}{2}H^2\Omm(a)\,\delta.
  \label{eq:growth_eq}
\end{equation}
%
The growing solution of Eq.~(\ref{eq:growth_eq}) defines the
\emph{linear growth factor} $\Dp(a)$, normalised to $\Dp = 1$ at a
chosen reference epoch:
%
\begin{equation}
  \ddot\Dp + 2H\dot\Dp
    = \tfrac{3}{2}H^2\Omm(a)\,\Dp.
  \label{eq:growth_D}
\end{equation}
%
In the matter-dominated limit $\Dp \propto a$.  The
\emph{logarithmic growth rate} is
%
\begin{equation}
  f(a) \equiv \frac{d\ln\Dp}{d\ln a} \approx \Omm(a)^{0.55}
  \quad\text{\citep{Linder2005}},
  \label{eq:fgrowth}
\end{equation}
%
where
%
\begin{equation}
  \Omm(a) = \frac{\Omm^0\,a^{-3}}{\Omm^0\,a^{-3}
    + \Omega_\Lambda + \Omega_k^0\,a^{-2}}.
  \label{eq:Omm_a}
\end{equation}
%
It is also convenient to define the \emph{conformal Hubble parameter}
$\calH(a) \equiv aH(a) = \dot{a}$, so that
$\dot\Dp = H f \Dp = (\calH/a) f \Dp$.

%----------------------------------------------------------------------
\section{Lagrangian Perturbation Theory Framework}
\label{sec:lpt}

\subsection{Displacement ansatz}

In comoving coordinates the position of a fluid element initially at
$\bq$ is
%
\begin{equation}
  \bx(\bq, t) = \bq + \bPsi(\bq, t),
  \label{eq:displacement}
\end{equation}
%
where $\bPsi$ is the (comoving) displacement field.  The peculiar
velocity is
%
\begin{equation}
  \bv_{\rm pec} = a\,\dot{\bx} = a\,\dot{\bPsi}.
  \label{eq:vpec_def}
\end{equation}

\subsection{Mass conservation and the Jacobian}

The density field follows from mass conservation.  Each fluid element
initially at $\bq$ carries fixed mass $\bar{\rho}\,d^3q$.  After
displacement it occupies the Eulerian volume element
%
\begin{equation}
  d^3x = \left|\det\!\left(\frac{\partial x_i}{\partial q_j}\right)\right|
          d^3q = |J|\,d^3q,
\end{equation}
%
where $J_{ij} = \partial x_i/\partial q_j$ is the deformation tensor.
Requiring $\rho(\bx)\,d^3x = \bar{\rho}\,d^3q$ gives
$\rho(\bx) = \bar{\rho}/|J|$, so
%
\begin{equation}
  1 + \delta(\bx) \equiv \frac{\rho(\bx)}{\bar{\rho}}
    = \left|\det\!\left(\delta_{ij} +
      \frac{\partial\Psi_i}{\partial q_j}\right)\right|^{-1}
  \equiv |J|^{-1},
  \label{eq:jacobian}
\end{equation}
%
which is \emph{exact} (no perturbative expansion has been made).

\subsection{Perturbative expansion}

In the matter-dominated limit the growing mode scales as $\Dp(t)
\propto a(t)$, and one expands the displacement as a series in
$\Dp$:
%
\begin{equation}
  \bPsi(\bq, t)
    = \Dp(t)\,\bPsi^{(1)}(\bq)
    + \Dpp(t)\,\bPsi^{(2)}(\bq)
    + \cdots,
  \label{eq:lpt_expansion}
\end{equation}
%
where $\Dpp \approx -\tfrac{3}{7}\Dp^2$ in $\Lambda$CDM (exact only in
matter domination; see \S\ref{sec:2lpt_growth}).

%----------------------------------------------------------------------
\section{First Order: the Zel'dovich Approximation}
\label{sec:za}

\subsection{Displacement field}

At first order, the displacement is the gradient of a scalar potential
$\phi^{(1)}$:
%
\begin{equation}
  \bPsi^{(1)}(\bq) = -\nabla_\bq\phi^{(1)}(\bq),
  \label{eq:psi1}
\end{equation}
%
with $\phi^{(1)}$ satisfying the Poisson equation in Lagrangian space:
%
\begin{equation}
  \nabla_\bq^2\phi^{(1)}(\bq) = \delta_{\rm lin}(\bq),
  \label{eq:poisson1}
\end{equation}
%
where $\delta_{\rm lin}$ is the linear density contrast at the
reference epoch at which $\Dp = 1$.
In Fourier space this is solved trivially:
%
\begin{equation}
  \boxed{
  \bPsi^{(1)}(\bk) = -\frac{i\bk}{k^2}\,\delta_{\rm lin}(\bk).
  }
  \label{eq:psi1_fourier}
\end{equation}
%
The IC generator draws $\delta_{\rm lin}(\bk)$ from a Gaussian with
variance $P_{\rm lin}(k)$ (the CLASS power spectrum), then applies
Eq.~(\ref{eq:psi1_fourier}) to compute particle displacements.

\subsection{Velocity field}

Differentiating Eq.~(\ref{eq:displacement}) with respect to time and
using $\bPsi = \Dp\bPsi^{(1)}$:
%
\begin{align}
  \bv_{\rm pec}^{(1)}
    &= a\,\dot{\Dp}\,\bPsi^{(1)}
     = a\,Hf\Dp\,\bPsi^{(1)}
     = \calH(a)\,f(a)\,\Dp\,\bPsi^{(1)},
  \label{eq:vel1}
\end{align}
%
where $\calH \equiv aH$ is the conformal Hubble parameter and $f$ is
the logarithmic growth rate from Eq.~(\ref{eq:fgrowth}).  At the IC
redshift this is the conversion MUSIC2 applies via
\texttt{cosmo\_vfact} $= \calH f / h$ to convert displacements to
velocities.

\subsection{Density field}

Linearising the Jacobian determinant in Eq.~(\ref{eq:jacobian}),
%
\begin{equation}
  \delta^{(1)}(\bx) \approx -\nabla_\bq\cdot\bPsi^{(1)}
  = \Dp\,\nabla_\bq^2\phi^{(1)}
  = \Dp\,\delta_{\rm lin}(\bq).
\end{equation}
%
At first order the density contrast is simply the linear field scaled
by the growth factor.

\subsection{Relation between the displacement potential and the
  gravitational potential}
\label{sec:phi_Phi}

The displacement potential $\phi^{(1)}$ satisfies
$\nabla^2\phi^{(1)} = \delta_{\rm lin}$.  On the other hand, the
linearised Poisson equation~(\ref{eq:poisson_lin}) gives
$\nabla^2\Phi = \tfrac{3}{2}H^2\Omm\,a^2\,\delta = \tfrac{3}{2}H^2\Omm\,a^2\Dp\,\delta_{\rm lin}$.
Comparing the two:
%
\begin{equation}
  \phi^{(1)}(\bq; z)
    = \frac{2}{3}\,\frac{\Dp(z)^{-1}}{H^2(z)\,\Omm(z)\,a^2}\,\Phi(\bq; z).
  \label{eq:phi_Phi}
\end{equation}
%
In practice, when $\Phi$ is evaluated at the same epoch as $\Dp$ (so
$\delta = \Dp\delta_{\rm lin}$ with the normalization
$\Phi = \frac{3}{2}H^2\Omm a^2\Dp\phi^{(1)}$), the factor $\Dp^{-1}$
cancels and the relation simplifies to
%
\begin{equation}
  \phi^{(1)}(\bq)
    = \frac{2}{3}\,\frac{1}{H^2(z)\,\Omm(z)\,a^2}\,\Phi(\bq; z),
  \label{eq:phi_Phi_simple}
\end{equation}
%
where $\phi^{(1)}(\bq)$ is the reference-epoch potential (at $\Dp=1$).
Substituting Eqs.~(\ref{eq:psi1}) and (\ref{eq:phi_Phi_simple}) into
Eq.~(\ref{eq:vel1}), the first-order peculiar velocity can be written
directly in terms of the proper gravitational potential:
%
\begin{equation}
  \bv_{\rm pec}^{(1)}(\bq; z)
    = -\frac{2}{3}\,\frac{f(z)}{a\,H(z)\,\Omm(z)}\,
      \nabla_\bq\Phi(\bq; z).
  \label{eq:vpec_Phi}
\end{equation}
%
This is the starting point for linear-theory peculiar-velocity
statistics.

\subsection{Shell crossing}

The Jacobian in Eq.~(\ref{eq:jacobian}) vanishes when the smallest
eigenvalue $\lambda_{\min}$ of the deformation tensor satisfies
$\Dp\,\lambda_{\min} = 1$, i.e.\ when two fluid elements with the
same $\bq$-plane coordinate arrive at the same Eulerian position.
This \emph{shell crossing} signals the breakdown of single-stream flow
and the onset of virialisation.  The ZA (being linear in $\bPsi$)
extends particles through the caustic without stopping them; higher
orders improve the accuracy but do not remove shell crossing.  For
this reason ICs must be laid down at sufficiently high redshift that
shell crossing has not yet occurred on any resolved scale.

%----------------------------------------------------------------------
\section{Second Order: 2LPT}
\label{sec:2lpt}

\subsection{Second-order displacement}

Expanding the equations of motion to second order yields an additional
irrotational displacement
%
\begin{equation}
  \bPsi^{(2)}(\bq) = -\nabla_\bq\phi^{(2)}(\bq),
  \label{eq:psi2}
\end{equation}
%
where the source of $\phi^{(2)}$ is quadratic in the first-order
potential:
%
\begin{equation}
  \boxed{
  \nabla_\bq^2\phi^{(2)}(\bq)
    = \sum_{i < j}
      \left[
        \frac{\partial^2\phi^{(1)}}{\partial q_i^2}
        \frac{\partial^2\phi^{(1)}}{\partial q_j^2}
        -
        \left(\frac{\partial^2\phi^{(1)}}{\partial q_i\,\partial q_j}\right)^{\!2}
      \right].
  }
  \label{eq:poisson2}
\end{equation}
%
The right-hand side is the sum of $2\times 2$ minors of the Hessian
of $\phi^{(1)}$, measuring the coupling between different components
of the tidal field.  In Fourier space this is a mode-coupling integral
(convolution), so $\bPsi^{(2)}$ mixes modes that are independent at
first order.

\subsection{Growth factor at second order}
\label{sec:2lpt_growth}

The second-order growing mode satisfies a separate ODE.  In the
matter-dominated ($\Omm = 1$) limit $\Dpp = -\tfrac{3}{7}\Dp^2$
exactly.  For $\Lambda$CDM a fitting formula accurate to
$\lesssim 0.3\%$ is \citep{Crocce2006}
%
\begin{equation}
  \Dpp(a) \approx -\frac{3}{7}\,\Dp^2(a)\,\Omm(a)^{-1/143}.
  \label{eq:D2}
\end{equation}
%
The corresponding logarithmic growth rate
$f^{(2)} \equiv d\ln|\Dpp|/d\ln a$ is well approximated by
\citep{Bouchet1995}
%
\begin{equation}
  f^{(2)}(a) \approx 2\,\Omm(a)^{4/7},
  \label{eq:f2}
\end{equation}
%
where $\Omm(a)$ is given by Eq.~(\ref{eq:Omm_a}).  In matter
domination $\Omm\to 1$ so $f^{(2)}\to 2$, exactly twice the
first-order rate $f\to 1$.  The 2LPT velocity correction is
%
\begin{equation}
  \bv_{\rm pec}^{(2)}
    = \calH(a)\,f^{(2)}(a)\,\Dpp\,\bPsi^{(2)}.
  \label{eq:vel2}
\end{equation}

\subsection{Total displacement and velocity}

Combining both orders:
%
\begin{align}
  \bx &= \bq + \Dp\,\bPsi^{(1)} + \Dpp\,\bPsi^{(2)},
  \label{eq:x_total}\\
  \bv_{\rm pec} &= \calH(a)\bigl(f\,\Dp\,\bPsi^{(1)}
                   + f^{(2)}\Dpp\,\bPsi^{(2)}\bigr).
  \label{eq:v_total}
\end{align}
%
The generalisation of Eq.~(\ref{eq:vpec_Phi}) to 2LPT is
%
\begin{equation}
  \bv_{\rm pec}(\bq; z)
    = -\frac{2}{3}\,\frac{1}{a\,H(z)\,\Omm(z)}\,
      \sum_{n=1}^{2} f^{(n)}(z)\,\nabla_\bq\Phi^{(n)}(\bq; z),
  \label{eq:vpec_2lpt}
\end{equation}
%
where $\Phi^{(n)}$ is the $n$-th-order proper gravitational potential
and $f^{(1)} \equiv f$.  In MUSIC2 the factor
$(6/7)/v_{\rm fac2lpt}$ in \texttt{main.cc} implements the ratio
$\Dpp/\Dp^2 = -3/7$ and the corresponding velocity coefficient.

%----------------------------------------------------------------------
\section{Relation to the Power Spectrum}
\label{sec:pk}

\subsection{What CLASS provides}

CLASS solves the linearised Boltzmann equations and outputs the
\emph{linear} matter power spectrum $P_{\rm lin}(k)$.  This is the
input to MUSIC2 and is the correct comparison for the 1LPT (ZA)
density field.

\subsection{Non-linear corrections from 2LPT}

Because the 2LPT displacement is quadratic in $\delta_{\rm lin}$, the
resulting density field contains contributions at second order in
$\delta_{\rm lin}$.  At one loop, the power spectrum of the displaced
particle distribution is \citep{Bernardeau2002}
%
\begin{equation}
  P_{\rm 2LPT}(k)
    = P_{\rm lin}(k)
    + P_{22}(k)
    + 2\,P_{13}(k)
    + \mathcal{O}(\delta_{\rm lin}^6),
  \label{eq:pk_loop}
\end{equation}
%
where the one-loop terms involve convolutions of $P_{\rm lin}$:
%
\begin{align}
  P_{22}(k) &= 2\int\!\frac{d^3q}{(2\pi)^3}\,
    \left[F_2(\bk-\mathbf{q},\,\mathbf{q})\right]^2
    P_{\rm lin}(|\bk-\mathbf{q}|)\,P_{\rm lin}(q), \\
  P_{13}(k) &= 3\,P_{\rm lin}(k)\int\!\frac{d^3q}{(2\pi)^3}\,
    F_3(\bk,\mathbf{q},-\mathbf{q})\,P_{\rm lin}(q),
\end{align}
%
with $F_2$, $F_3$ the standard perturbation theory (SPT) mode-coupling
kernels.  Consequently the power spectrum measured from the IC
particles will lie \emph{above} the CLASS linear theory at high $k$,
even before any non-linear evolution.

\subsection{Size of the correction}

The one-loop terms scale as $P_{\rm lin}^2 \propto \sigma_8^4(z)$.
Since $\sigma_8(z) = \sigma_8(0)\,\Dp(z)$ and $\Dp(z) \approx
1/(1+z)$ in matter domination:
%
\begin{equation}
  \frac{P_{22}+2P_{13}}{P_{\rm lin}} \sim \sigma_8^2(z)
    = \left[\frac{\sigma_8(0)}{1+z}\right]^2.
\end{equation}
%
For our CV\_22 cosmology ($\sigma_8 = 0.8$):
%
\begin{center}
\begin{tabular}{lcc}
\hline
Redshift & $\sigma_8(z)$ & One-loop correction \\
\hline
$z = 127$ & 0.006 & $\lesssim 0.004\%$ \\
$z = 45$  & 0.017 & $\lesssim 0.03\%$  \\
$z = 2$   & 0.27  & $\sim 7\%$ (high $k$) \\
$z = 0.5$ & 0.55  & $\sim 30\%$ (high $k$) \\
\hline
\end{tabular}
\end{center}
%
For $z \gtrsim 10$ the linear approximation is excellent.  At $z = 2$
(our default pipeline redshift) deviations are visible at $k \gtrsim
0.5\,h\,\mathrm{Mpc}^{-1}$.  A consistent comparison at $z \lesssim 5$
requires a one-loop prediction from a code such as \textsc{PyBird},
\textsc{velocileptors}, or CLASS with \texttt{halofit}.

%----------------------------------------------------------------------
\section{Why 2LPT Matters for ICs}

Even when the one-loop correction to $P(k)$ is small, starting
simulations with ZA-only ICs introduces \emph{transients}: spurious
mode-coupling terms that take several dynamical times to die away
\citep{Crocce2006}.  These arise because ZA misplaces particles at
second order, and the $N$-body force then has to ``correct'' this
error.  2LPT removes the leading transient and is standard practice
for $z_{\rm start} \lesssim 100$.  For $z_{\rm start} \gtrsim 100$
the growth factor is small enough that ZA suffices.

%----------------------------------------------------------------------
\section{Summary}

\begin{itemize}
  \item \textbf{Fluid equations}: the linearised continuity, Euler, and
        Poisson equations in an expanding background lead to the growth
        equation~(\ref{eq:growth_D}).  The growing solution $\Dp(a)$
        has logarithmic growth rate $f \approx \Omm(a)^{0.55}$
        \citep{Linder2005}.

  \item \textbf{1LPT / ZA}: $\bPsi^{(1)} = -i(\bk/k^2)\delta_{\rm lin}(\bk)$,
        velocity $\bv_{\rm pec}^{(1)} = \calH f\Dp\bPsi^{(1)}$.
        The displacement potential is related to the proper gravitational
        potential by $\phi^{(1)} = (2/3)(H^2\Omm a^2)^{-1}\Phi$,
        giving $\bv_{\rm pec}^{(1)} = -(2/3)\,f/(aH\Omm)\,\nabla\Phi$.
        Exact for planar collapse; breaks down at shell crossing.

  \item \textbf{2LPT}: adds $\bPsi^{(2)} = -\nabla\phi^{(2)}$ with source
        Eq.~(\ref{eq:poisson2}).  Growth mode
        $\Dpp \approx -\tfrac{3}{7}\Dp^2\,\Omm^{-1/143}$
        \citep{Crocce2006}, growth rate $f^{(2)} \approx 2\Omm(a)^{4/7}$
        \citep{Bouchet1995}, velocity
        $\bv_{\rm pec}^{(2)} = \calH f^{(2)}\Dpp\bPsi^{(2)}$.
        Eliminates leading transients and is correct through one loop
        in $P(k)$.

  \item \textbf{CLASS vs.\ MUSIC2}: CLASS outputs linear $P_{\rm lin}(k)$
        (first order in $\delta$); the IC particle distribution has
        additional power $\sim\sigma_8^2(z)$ from 2LPT mode coupling.
        The discrepancy is negligible at $z \gtrsim 10$ and a few percent
        at $z = 2$.

  \item \textbf{SWIFT velocity storage}: SWIFT ICs store
        $\bv_{\rm int} = a\,\bv_{\rm pec}$, not the peculiar velocity
        directly.  The velocity correlation function $\psi(r) =
        \langle\bv_{\rm pec}\cdot\bv_{\rm pec}'\rangle$ must therefore
        be extracted by dividing the stored correlation by $a^2$.
\end{itemize}

%----------------------------------------------------------------------
\bibliography{lpt_review}

\end{document}
